\section{Related Work}
\label{sec:relwork}
This work draws on edge IDS, conditional inference, model serving, and learned resource orchestration. Recent IDS research addresses continual detection, heterogeneous representations, and runtime model selection~\cite{Wang2024BFSNID,AOCIDS2024,StatAvg2025,DMIDS2025}. Auto-IDaS and FEDMS show that learned or confidence-driven detector selection is operationally plausible~\cite{Lin2024AutoIDaS,Hasan2026FEDMS}. However, neither controller complexity nor a lower-cost operating point alone establishes the value of adaptation. We therefore compare each selector with the static detector it actually uses and report the full state–action map alongside the occupancy of visited states.

Conditional-computation systems such as BranchyNet, adaptive network selection, MSDNet, and shallow–deep networks escalate uncertain inputs to costlier models~\cite{Teerapittayanon2016BranchyNet,Bolukbasi2017Adaptive,Huang2018MSDNet,Kaya2019ShallowDeep}. Our fixed confidence cascade applies this approach to whole-device IDS measurement. It also reveals a cost that row-level escalation rates can obscure. Only 1.34\% of rows are escalated, but these rows activate MedRF in 66.76\% of windows. Thus, call granularity, rather than the row fraction alone, governs the observed cost.

Serving systems such as Clipper, INFaaS, and Clockwork explicitly define service-level objectives (SLOs) for accuracy, latency, and cost, and route requests among model variants~\cite{Crankshaw2017Clipper,Romero2021INFaaS,Gujarati2020Clockwork}. We adopt this service-contract perspective and add two controls missing from these systems: a matched static action to separate the value of routing from model choice, and a pre-deployment audit to determine whether the declared reward allows preference reversals on validation support.

RL-based edge orchestration includes virtualized RAN (vRAN) reconfiguration, vehicle--edge scheduling, dynamic voltage and frequency scaling (DVFS), offloading, and QoS control~\cite{Murti2024vRAN,Zhang2024VehicleEdge,Zhang2024DVFO,Lou2024Startup,Kim2024DistributedOffload,Maleki2024QCS,Nan2026RobustDNN,Wang2025TrustOffload,Hu2025SITOff}. Tabular Q-learning and DQN are representative discrete and neural value-based controllers~\cite{Watkins1992,Mnih2015}. Our contribution to this line of work is a diagnostic methodology combining an analytic reward test, matched physical comparisons, a load–prevalence factorial design fixed before execution with an explicit deadline SLO, and an audit of actions outside training support. These components address known weaknesses in the evaluation of security ML and small-sample RL studies~\cite{Arp2022,Agarwal2021}.
